\subsection{Scene III --- The Editable City}\label{subsec:scene_city}

The third scene, \texttt{City\_Sandbox}, is a sandbox in which the user moves
through a city, rearranges it, and traverses it as an avatar, choosing between
these through the wrist menu of~\autoref{subsec:interaction_layer}.
\texttt{BraceletMenuCity} binds the three menu options --- traverse, move, and
grappling --- and falls back to a plain orbiting camera when none is active.

The traverse mode turns the hands into a flight stick. \texttt{CameraRingController}
reads both hands: with both closed into fists the player rig rotates --- yaw from
the angle of the line joining the hands, pitch from the vertical position of their
midpoint, roll from the line's tilt --- and with a single fist it translates through
the scene. Deadzones on each axis absorb small tracking movements, so the view holds
steady when the user means it to.

The move mode enables the shared pointer together with \texttt{BuildingGrabber}.
Confirming with \texttt{Victory} on a building picks it up; it then follows the hand,
in whichever of the pointer's two styles is active --- tracking the palm directly,
or carried by the displacement of the laser's hit point --- with its height held
fixed so it slides across the ground plane rather than lifting off. The motion is
smoothed, and releasing the confirmation sets the building down.

\begin{figure}[H]
	\centering
	\captionsetup{justification=centering}
	\includegraphics[width=0.85\linewidth, frame]{img/scene_city_move.png}
	\caption{A building selected and repositioned
	         across the ground plane in move mode.}
	\label{fig:city-move}
\end{figure}

The grappling mode equips the same bear character with omni-directional mobility
gear, and is built on the grapple core of Scene~I through a separate subclass,
\texttt{BearODMController}. Where the first scene threw one rope at a time, this
mode fires two together on both fists and winches the avatar toward the wall by
steadily shortening the ropes. Once both are anchored it tracks the avatar against
the wall, and when the swing carries it past, both hooks release with a forward
launch and time slows to a fraction of real speed for an aerial moment in which a
\texttt{Pointing\_Up} on either hand steers the body; a new throw is withheld until
the hands open again. The shared base supplies the ropes, the aim assist, and the
collision haptics, while the subclass supplies only the dual-hook winching, the
launch, and the slow-motion steering --- an arrangement that lets two quite
different traversal idioms reuse one body of grappling code.

This scene builds on Scene~I in two ways. First, it introduces a richer set of
hand-driven movements: paired grapples, winching, and a mid-air-steered launch,
all controlled by the single webcam
that~\autoref{subsubsec:sota_immersive} compares against motion controllers and
headset rigs. Second, objects in the scene can be picked up and repositioned
directly with the hands --- the same controller-free manipulation demonstrated in
the deformer, but applied to whole objects rather than individual mesh vertices.

\begin{figure}[H]
	\centering
	\captionsetup{justification=centering}
	\includegraphics[width=0.8\linewidth, frame]{img/scene_city_odm.png}
	\caption{Omni-directional mobility: the avatar winches
	         toward a wall on twin ropes.}
	\label{fig:city-odm}
\end{figure}
